\chapter{Data and Preprocessing}
\label{ch:data}

\section{The MailEx corpus}
All data derive from MailEx \citep{srivastava2023mailex}, built on Enron email \citep{klimt2004enron}. Our conversion recomputes exactly the published splits: 3,117/414/405 emails in 1,200/150/150 threads, of which 636/70/70 carry no event. Each email becomes a unified JSONL record with tokens, character offsets, thread context and the original event labels. Two properties shaped later design: no record has a sending timestamp, so relative dates in the corpus cannot be resolved; and sender/recipient fields are noisy raw headers. A conservative cleaner removes disclaimers, quoted blocks and signatures without ever deleting an annotated token; all 3,936 records were re-aligned with no fallbacks.


\begin{figure}[htbp]
\centering
\chart[0.85\linewidth]{mailex_event_types}
\caption{Emails containing each MailEx event type, per split (recomputed from the converted files).}
\label{fig:events}
\end{figure}
\section{Three label sets}
\begin{table}[htbp]
\centering
\caption{Label sets. Major-error rates are from AI-assisted inspections of 30 to 40 emails and are indicative only.}
\label{tab:labelsets}
\footnotesize
\begin{tabularx}{\linewidth}{@{}p{24mm}p{24mm}p{24mm}Lp{14mm}@{}}
\toprule
Set & Emails (tr/dev/test) & Types & Origin & Major errors \\
\midrule
SILVER-v1 & 450 / 60 / 90 & 6 entity, 4 relation & automatic + LLM annotation of 600 selected emails, schema validation & 73\% \\
SILVER-v2 & 449 / -- / -- & TASK, PERSON, TEMPORAL & repair of v1 (audits, validated LLM output, filtering) & 33\% \\
MailEx-derived & 2,845 / 389 / -- & TASK, PERSON, TEMPORAL & deterministic mapping of MailEx human arguments & ${\approx}9\%$ \\
\bottomrule
\end{tabularx}
\end{table}

\Cref{tab:labelsets} summarises the label sets. SILVER-v1 came first; its first model failed, a repair produced SILVER-v2, which trained the entity checkpoint the application uses, and the relation dataset was built from SILVER-v1. Because SILVER labels remained noisy, a second lineage maps MailEx's human annotations directly (mapping v1.2.0): TASK from \lab{Action Description} in request events plus commitments starting with \emph{will}, \emph{'ll} or \emph{going to}; TEMPORAL from \lab{Date}/\lab{Time} slots, dropping event clauses (\emph{once \ldots}) and spans over 8 tokens; PERSON from \lab{Members} slots that pass a capitalised-name filter. In training this kept 1,664 TASK, 1,040 PERSON and 1,036 TEMPORAL spans and rejected 3,349 non-name members and 2,101 non-commitment deliveries. The resulting labels are cleaner but incomplete: about one email in eight contains an unannotated request, and names and dates are labelled only when they are event arguments.


\begin{figure}[htbp]
\centering
\diagram{dataset_pipeline}
\caption{Dataset preparation. Lineage A produced SILVER-v1, its repair SILVER-v2 and the relation data; Lineage B maps MailEx human arguments directly and fed the Phase 2 experiments. The leakage audit was added after the first results.}
\label{fig:datapipeline}
\end{figure}

\begin{table}[htbp]
\centering
\caption{MailEx-to-Text2Task mapping, version 1.2.0.}
\label{tab:mapping}
\footnotesize
\begin{tabularx}{\linewidth}{@{}p{16mm}LL@{}}
\toprule
Type & Source in MailEx & Filters \\
\midrule
TASK & \lab{Action Description} in request events; deliveries starting with a commitment form (\emph{will}, \emph{'ll}, \emph{shall}, \emph{going to}) & excludes \emph{will be \textless state\textgreater}, \emph{will have}, \emph{will not}; a requests-only variant is the ablation \\
TEMPORAL & any \lab{Date} or \lab{Time} slot & drop if over 8 tokens or a subordinate event clause without a calendar or clock word \\
PERSON & any \lab{Members} slot & 1--3 capitalised tokens; rejects pronouns, roles, company/place words, acronyms; repeated names in the email also tagged \\
\bottomrule
\end{tabularx}
\end{table}

\begin{figure}[htbp]
\centering
\chart[0.92\linewidth]{entity_distribution}
\caption{Entity span counts in the three label sets. GROUP and EVENT are too rare to learn.}
\label{fig:entitydist}
\end{figure}
\section{Tokenisation and BIO tagging}
Texts are tokenised with the uncased DistilBERT WordPiece vocabulary and offset mapping. The first piece of each span is labelled \lab{B-TYPE}, later pieces \lab{I-TYPE}, everything else \lab{O}; special tokens get the ignore index. Texts longer than 512 pieces are windowed. A reconstruction check decodes labels back to character spans and found 0 mismatches over all 3,262 MailEx-derived rows. \Cref{fig:bio} follows one real training email from MailEx slots to BIO labels: the commitment \emph{'ll cover this} becomes a TASK, \emph{tomorrow} a TEMPORAL, \emph{Andrea} a PERSON, while the pronoun members \emph{We} and \emph{I} are rejected.

\begin{figure}[htbp]
\centering
\diagram{bio_tagging}
\caption{From MailEx arguments to WordPiece BIO labels for a real training record (TASK blue, PERSON teal, TEMPORAL ochre).}
\label{fig:bio}
\end{figure}

\section{Relations, imbalance and leakage}
For relations, every TASK$\to$PERSON and TASK$\to$TEMPORAL pair in an email is a candidate, marked with typed tokens (\code{[E1\_TASK]} \ldots{} \code{[/E1\_TASK]}, \code{[E2\_TEMPORAL]} \ldots) and labelled with its relation or \lab{NONE}. Training kept all 1,144 positives and 1,144 sampled negatives; \lab{DUE\_ON} has only 33 training pairs and \lab{DEPENDS\_ON} only 4, which is why dependencies are rule-based. At the token level 91.0\% of 245,834 labels are \lab{O} and each \lab{B-} label is under 0.7\% (\cref{fig:imbalance}). Finally, an audit after the first results found that 39 of 389 development emails copied training emails; all MailEx-derived scores below use the 349 clean emails.

\begin{figure}[htbp]
\centering
\begin{subfigure}[b]{0.49\linewidth}\centering\chart{bio_imbalance}\caption{BIO labels (MailEx-derived train, log scale)}\end{subfigure}\hfill
\begin{subfigure}[b]{0.49\linewidth}\centering\chart{relation_distribution}\caption{Relation candidates by label (log scale)}\end{subfigure}
\caption{Class imbalance in the entity and relation data.}
\label{fig:imbalance}
\end{figure}
